\documentclass[../main.tex]{subfiles}
\graphicspath{{\subfix{../Images/}}}

\begin{document}
\onecolumn
% =================================================================
% Section IV -- Problem Formulation (rewritten, REVISION 2)
%
% Companion to method_implementation_skeleton.tex / Method_and_Implementation.md.
%
% This file is a concrete draft of fixes to problem_formulation.tex. It is a
% full drop-in replacement, not a diff: every label from the original is kept
% (so the Method skeleton's \eqref/\ref calls do not break), and a small number
% of new labels are added for objects that previously had no formal equation.
% See GAPS.tex for the rationale behind every change below; each change is
% tagged inline with [Gxx] (internal-consistency fix) or [Kxx] (fix motivated
% by comparison against Knippenberg, Holenderski & Menkovski 2021, the
% non-lattice-MAPF-with-DRL paper this section's environment/observation model
% is adapted from) matching the numbered entry in GAPS.tex.
%
% REQUIRES A BIBLIOGRAPHY ENTRY: \citet{knippenberg2021} is used three times
% below (environment connectivity, scope/deadlines x2), in natbib author-year
% form. The original problem_formulation.tex has no citations at all, so this
% file cannot confirm what citation package main.tex loads; if it is not
% natbib, replace \citet{knippenberg2021} with \cite{knippenberg2021} (or this
% thesis's own citation macro) at each site. Add to your .bib:
%   @inproceedings{knippenberg2021,
%     author    = {van Knippenberg, Marijn and Holenderski, Mike and Menkovski, Vlado},
%     title     = {Time-Constrained Multi-Agent Path Finding in Non-Lattice
%                   Graphs with Deep Reinforcement Learning},
%     booktitle = {Proceedings of the 13th Asian Conference on Machine
%                   Learning (ACML)},
%     series    = {Proceedings of Machine Learning Research},
%     volume    = {157},
%     year      = {2021}
%   }
%
% CONTRACT (unchanged from the original)
%   This section says WHAT the objects are, not HOW anything is computed.
%   Every "how" is a forward reference to Section V (Method) or Section VI
%   (Implementation). Per Method_and_Implementation.md, a full multi-entry
%   notation table is Appendix material, not Problem Formulation material --
%   this file only points to Appendix~\ref{app:notation} (see GAPS.tex for the
%   appendix-ready table itself, since this repo has no appendix.tex to edit).
%
% Labels exported by this section (unchanged ones first, new ones marked NEW):
%   eq:onestepcost  eq:actions      eq:transition    eq:vertexconflict
%   eq:swapconflict eq:feasiblestorage  eq:storageupdate  eq:storageobjective
%   eq:task         eq:coupling     eq:servicetime   eq:partition
%   eq:localsubgraph eq:observation eq:potential     eq:commgraph
%   eq:posg         eq:congestion   eq:objective     eq:throughput
%   eq:movementcost eq:entropy      eq:concentration eq:functional
%   eq:decisionproblem  eq:rqformal
%   eq:lifecycle (NEW, [G5])  eq:waiting (NEW, [G7])  eq:potentialbound (NEW)
% Section labels: sec:problem, sec:pf:env, sec:pf:agents, sec:pf:storage,
%   sec:pf:tasks, sec:pf:controllers, sec:pf:observations, sec:pf:game,
%   sec:pf:coupling, sec:pf:measures, sec:pf:notation (NEW, now a pointer to
%   Appendix~\ref{app:notation} rather than an inline table -- see [G11]/GAPS.tex),
%   sec:pf:scope
%
% Requires: amsmath, amssymb, graphicx, booktabs
% =================================================================

\section{Problem Formulation}
\label{sec:problem}

This section defines the coupled storage--routing system studied in this thesis.
It proceeds in the order in which a simulation run unfolds: the environment, the
agents that move in it, the storage state, the work that arrives, the controllers
that decide, what a decentralised agent can see, the resulting learning problem,
the loop that closes between the two layers, and the quantities used to score a
run. Each object is described in plain terms before it is written formally.

Throughout, this section states \emph{what} the objects are, not \emph{how} they
are computed. Realisations --- heuristics, algorithms, network architectures and
hyperparameters --- are deferred to Section~\ref{sec:method}, and software
concerns to Section~\ref{sec:implementation}.

% -----------------------------------------------------------------
\subsection{The environment}
\label{sec:pf:env}

The warehouse is a finite \emph{directed} graph $G=(V,E)$. % [G10]
For a one-way aisle, $(v,w) \in E$ does not imply $(w,v) \in E$, and the two
directions may carry different costs; nothing below assumes $E$ is symmetric.
Vertices are the discrete positions a robot can occupy: aisle cells, cross-aisle
junctions, storage faces and drop-off points. Edges are the movements that are
physically possible between two adjacent positions. A rectangular grid is one
special case, but the formulation deliberately does not assume a lattice,
because real warehouses contain one-way aisles, diagonal shortcuts and irregular
blocks.

Four roles label the vertices by what can happen at them; unlike a partition, % [G1]
these roles are not required to be pairwise disjoint, and this section no longer
calls them one. $V_{\mathrm{mov}} \subseteq V$ are the vertices an agent may
occupy. $V_{\mathrm{str}} \subseteq V_{\mathrm{mov}}$ are storage faces, where
items are held and from which pickups are made: a pickup at $s_j$ is completed
by an agent occupying $s_j$, so storage faces must themselves be occupiable, and
$V_{\mathrm{str}}$ is stated as a subset of $V_{\mathrm{mov}}$ rather than merely
of $V$. $V_{\mathrm{del}} \subseteq V_{\mathrm{mov}}$ are delivery or drop-off
points, where completed tasks terminate, occupiable for the same reason.
$V_{\mathrm{ep}} \subseteq V_{\mathrm{mov}}$ are endpoints: vertices at which an
idle agent may remain without blocking transit. $V_{\mathrm{str}}$,
$V_{\mathrm{del}}$ and $V_{\mathrm{ep}}$ need not be pairwise disjoint from one
another --- a generated instance will typically keep them so, but nothing in the
formulation below relies on it. Each edge $e \in E$ carries a traversal cost
$c(e) > 0$, and $d_G(v,w)$ denotes the shortest-path cost from $v$ to $w$ along
directed edges of $G$. Throughout, $G$ is assumed rich enough that $d_G(v,w)$ is
finite for every pair of vertices that appears in the constructions below (in
particular, every occupiable vertex can reach every vertex it is ever assigned
as a target); generated instances are checked against this in
Section~\ref{sec:impl:instances}. % [G10]
This is deliberately weaker than the blanket strong-connectivity assumption used
for the same purpose by \citet{knippenberg2021}, whose non-lattice MAPF
formulation this section otherwise closely follows for the environment and
observation model: strong connectivity would be simpler to state and verify, at
the cost of ruling out layouts with an unreachable side-room or a
maintenance-only spur, which the ``irregular blocks'' motivation above is meant
to admit. If that generality is not actually exercised by the generated
instances, adopting strong connectivity outright is the simpler and equally
sufficient choice and is worth reconsidering in Section~\ref{sec:impl:instances}. % [K1]

Figure~\ref{fig:instance} shows a small instance with all four roles marked.

\begin{figure}[t]
  \centering
  \includegraphics[width=\linewidth]{Images/instance.png}
  \caption{A small warehouse instance. Vertex shape encodes role: storage faces,
  delivery points, endpoints and plain transit vertices. A single task
  $\tau_j = (r_j, s_j, g_j, k_j)$ is shown as a route from its pickup vertex
  $s_j$ to its delivery vertex $g_j$. The layout is not a lattice: aisle
  structure, not grid geometry, determines which movements exist.}
  \label{fig:instance}
\end{figure}

At each timestep an agent at vertex $v$ either waits or traverses one outgoing
edge. Writing $N^{+}(v) = \{w \in V_{\mathrm{mov}} : (v,w) \in E\}$ for the
movable successors of $v$, the action set at $v$ is
\begin{equation}
  \mathcal{U}(v) = \{\mathrm{wait}\} \cup \{\mathrm{move}(w) : w \in N^{+}(v)\}.
  \label{eq:actions}
\end{equation}
Waiting costs $c_{\mathrm{wait}} \geq 0$ and moving from $v$ to $w$ costs
$c(v,w)$, so the realised one-step cost is
\begin{equation}
  \hat{c}(v,w) =
  \begin{cases}
    c_{\mathrm{wait}}, & v = w, \\
    c(v,w),            & (v,w) \in E.
  \end{cases}
  \label{eq:onestepcost}
\end{equation}
The two cases are needed because a self-loop $(v,v)$ is not an element of $E$, so
$c$ is undefined on a wait and any route cost containing a wait would otherwise be
ill-typed.\footnote{The value of $c_{\mathrm{wait}}$ relative to edge costs
determines whether an agent prefers to yield or to detour; it is therefore a
parameter of the study rather than a constant of the model, and is reported in
Section~\ref{sec:method:training}.}

Some baseline algorithms are complete only on \emph{well-formed} instances. Where
such a baseline is used, the following is imposed as an assumption on the
generated instance rather than derived: at least $m$ non-task endpoints exist
($m = |\mathcal{A}|$ is the fleet size, introduced in
Section~\ref{sec:pf:agents}), and between any pair of endpoints there is a path
that does not pass through a third endpoint. Instances are checked against this
condition in Section~\ref{sec:impl:instances}, and the condition is reported only
for the baselines to which it applies.

% -----------------------------------------------------------------
\subsection{Agents and collisions}
\label{sec:pf:agents}

A fleet $\mathcal{A} = \{a_1,\dots,a_m\}$ of homogeneous agents moves on $G$.
Agent $a_i$ occupies $\ell_i(t) \in V_{\mathrm{mov}}$ at time $t$ and selects
$u_i(t) \in \mathcal{U}(\ell_i(t))$. Its realised route over the horizon is
$p_i = (\ell_i(0), \dots, \ell_i(T))$. A single agent's transition is legal when
\begin{equation}
  \ell_i(t+1) = \ell_i(t)
  \quad \vee \quad
  \bigl(\ell_i(t), \ell_i(t+1)\bigr) \in E .
  \label{eq:transition}
\end{equation}

Legality of each agent's own move is not sufficient, because agents share the
infrastructure. Two agents collide if they end a timestep at the same vertex, or
if they exchange positions across a single edge. A joint transition is therefore
\emph{collision-free} when both
\begin{align}
  \ell_i(t+1) &\neq \ell_j(t+1) && \forall i \neq j,
  \label{eq:vertexconflict} \\
  \neg\bigl( \ell_i(t) = \ell_j(t+1) \;\wedge\; \ell_j(t) = \ell_i(t+1) \bigr)
      & && \forall i \neq j .
  \label{eq:swapconflict}
\end{align}
The two conditions are logically independent. Constraint
\eqref{eq:vertexconflict} inspects the successor state alone;
\eqref{eq:swapconflict} inspects the pair (state, successor) and is invisible to
any check performed only at $t+1$. Omitting it permits two agents to pass through
one another in a single-lane aisle. Figure~\ref{fig:conflicts} illustrates both,
together with the feasible alternative in which one agent waits.

\begin{figure}[t]
  \centering
  \includegraphics[width=\linewidth]{Images/conflicts.png}
  \caption{The two forbidden joint transitions and one feasible alternative.
  Panel~(c) shows the resolution the model always admits: waiting is a legal
  action at every vertex, so a collision-free joint action always exists. Which
  agent yields is a decision, not a constraint, and is fixed in
  Section~\ref{sec:method:resolution}.}
  \label{fig:conflicts}
\end{figure}

Constraints \eqref{eq:vertexconflict}--\eqref{eq:swapconflict} are hard: a run in
which they are violated is not a solution to this problem. They are stated here
as feasibility conditions on the realised trace. A learned policy induces a
distribution over joint actions whose support may contain infeasible ones, so
feasibility cannot be delegated to a penalty term in a reward; the mechanism that
enforces it is specified in Sections~\ref{sec:method:rl}
and~\ref{sec:method:resolution} and verified in Section~\ref{sec:impl:validation}.
That resolution mechanism is deterministic given the state (Section~\ref{sec:pf:scope}), and it is applied identically regardless of which $\pi \in \Pi$
is in force; this fact is used in Section~\ref{sec:pf:measures} to clarify what
work \eqref{eq:vertexconflict}--\eqref{eq:swapconflict} do inside the formal
decision problem.

% -----------------------------------------------------------------
\subsection{Where the items are}
\label{sec:pf:storage}

Items are held at storage faces, and where they are held determines where pickups
happen. Let $\mathcal{K}$ be the set of SKU types. The storage state must be able
to represent split inventory, multiple units, and finite shelf space, so it is
defined as a counting function
\[
  x_t : \mathcal{K} \times V_{\mathrm{str}} \to \mathbb{N}_0 ,
\]
where $x_t(k,v)$ is the number of units of SKU $k$ held at storage vertex $v$ at
time $t$. Let $b_k > 0$ be the capacity consumed by one unit of SKU $k$ and
$\mathrm{cap}(v) > 0$ the capacity of vertex $v$. The set of feasible
configurations is
\begin{equation}
  \mathcal{X} =
  \Bigl\{ x \in \mathbb{N}_0^{|\mathcal{K}| \times |V_{\mathrm{str}}|} :
    \sum_{k \in \mathcal{K}} b_k\, x(k,v) \le \mathrm{cap}(v)
    \quad \forall v \in V_{\mathrm{str}} \Bigr\}.
  \label{eq:feasiblestorage}
\end{equation}
A simpler mapping $\sigma_t : \mathcal{K} \to V_{\mathrm{str}}$ would force one
location per SKU, unit quantity and unlimited shelf space; it is admissible only
if those three restrictions are stated explicitly. Note that
\eqref{eq:feasiblestorage} is a system of $|V_{\mathrm{str}}|$ knapsack
constraints, so capacity couples the placement decisions of different SKUs at a
shared vertex. The system is initialised with a given feasible configuration % [G4]
$x_0 \in \mathcal{X}$, supplied as part of the problem instance; the update rule
below evolves $x_0$ forward at each storage epoch.

The storage layer acts on a slow clock. At epochs
$t \in \{\Delta, 2\Delta, \dots\}$ the configuration is updated by a rule $F$,
\begin{equation}
  x_t = F\bigl(x_{t-\Delta},\, \hat{\rho}_t,\, \hat{\mu}_t,\, \hat{w}_t\bigr)
        \in \mathcal{X} ,
  \label{eq:storageupdate}
\end{equation}
where $\hat{\rho}_t$ estimates per-SKU demand, $\hat{\mu}_t(e)$ estimates
traversal counts per edge, and $\hat{w}_t(e)$ estimates waiting attributed to an
edge or its downstream contention point. All three are computed only from data
available strictly before the epoch, so \eqref{eq:storageupdate} is causal. The
fixed-storage baseline is the identity rule $F_{\mathrm{fix}}(x,\cdot) = x$,
defined directly rather than as the informal limit $\Delta = \infty$. Formally, % [G3]
$\mathcal{F}$ is the set of measurable functions
$F : \mathcal{X} \times \mathbb{R}_{\ge 0}^{|\mathcal{K}|} \times
\mathbb{R}_{\ge 0}^{|E|} \times \mathbb{R}_{\ge 0}^{|E|} \to \mathcal{X}$ taking a
prior configuration and the three estimates to a feasible configuration;
$F_{\mathrm{fix}} \in \mathcal{F}$ is one particular element, and the decision
problem of Section~\ref{sec:pf:measures} ranges over all of $\mathcal{F}$.

Because capacity couples SKUs, the update cannot be a per-item greedy score
applied independently. It is a joint constrained decision over $\mathcal{X}$,
\begin{equation}
  x_t \in \arg\min_{x \in \mathcal{X}}
    \bigl[\, \alpha D(x; \hat{\rho}_t)
       + \beta K(x; \hat{\mu}_t, \hat{w}_t)
       + \chi R(x, x_{t-\Delta}) \,\bigr],
  \label{eq:storageobjective}
\end{equation}
where $D$ measures demand-weighted access distance, $K$ measures predicted
congestion exposure, and $R$ penalises reassignment. In words: $D$ says popular
items should be close, $K$ says they should not all sit on the same corridor, and
$R$ says the warehouse should not be rearranged every epoch. Because
\eqref{eq:storageobjective} is an $\arg\min$, it may admit more than one % [G3]
minimiser; $F$ is required to select a single element of the minimising set via a
fixed tie-breaking rule, so that \eqref{eq:storageupdate} defines a genuine
function rather than a set-valued correspondence. The three functional forms, the
weights $\alpha, \beta, \chi$, the tie-breaking rule, and the heuristic that
actually searches $\mathcal{X}$ are specified in Section~\ref{sec:method:storage};
\eqref{eq:storageobjective} is an integer program and is not solved to
optimality. Setting $\chi = 0$ removes relocation cost from the model entirely
and is recorded in Section~\ref{sec:pf:scope} as a limitation rather than a
simplification in passing.

% -----------------------------------------------------------------
\subsection{Where the work comes from}
\label{sec:pf:tasks}

Orders arrive continuously. At each timestep, orders are drawn from a
distribution $P_{\mathrm{ord}}$, an order being a finite multiset of
SKU--quantity pairs, and each order is converted into one or more transport tasks
by a task-generation rule. A task is a quadruple % [G2]
\begin{equation}
  \tau_j = (r_j,\, s_j,\, g_j,\, k_j), \qquad
  r_j \in \mathbb{N}_0,\quad s_j \in V_{\mathrm{str}},\quad
  g_j \in V_{\mathrm{del}}, \quad k_j \in \mathcal{K},
  \label{eq:task}
\end{equation}
consisting of a release time, a pickup vertex, a delivery vertex, and the
requested SKU $k_j$. The expected number of tasks released per timestep is
$\lambda_{\mathrm{task}}$.

The pickup vertex is not free. For the requested SKU $k_j$ it must satisfy
\begin{equation}
  x_{r_j}(k_j,\, s_j) > 0 .
  \label{eq:coupling}
\end{equation}
Equation~\eqref{eq:coupling} is the formal link between the two layers and the
central object of this thesis: it makes the spatial distribution of pickup
vertices a function of the storage configuration, so any change to $x_t$ changes
the routing problem that the controller subsequently faces. If inventory
depletion and replenishment are not simulated, $x_t$ is interpreted as a
placement template rather than a physical inventory level, and this
interpretation is stated with the results.

A task passes through three disjoint sets over its life, defined directly from % [G5]
three timestamps rather than only narratively. Let $y_j$ be the assignment time
of $\tau_j$ --- the timestep at which task assignment selects an agent for it ---
and $d_j$ the completion (delivery) time, so that $r_j \le y_j \le d_j$. Then
\begin{equation}
  \mathcal{Q}_t = \{\tau_j : r_j \le t < y_j\}, \qquad
  \mathcal{B}_t = \{\tau_j : y_j \le t < d_j\}, \qquad
  \mathcal{C}_t = \{\tau_j : d_j \le t\},
  \label{eq:lifecycle}
\end{equation}
the waiting, active and completed sets. These are pairwise disjoint by
construction, since the three defining half-open intervals partition
$[r_j, \infty)$ for every task --- disjointness is therefore a consequence of
\eqref{eq:lifecycle}, not a separately asserted property. The convention used
throughout is that execution begins at \emph{assignment} ($t = y_j$), not at
arrival of the agent at $s_j$. Writing $d_j$ for the completion time as above,
the service time is
\begin{equation}
  \zeta_j = d_j - r_j ,
  \label{eq:servicetime}
\end{equation}
which spans both phases and therefore cannot be improved by delaying assignment.
Figure~\ref{fig:lifecycle} shows the lifecycle and the interval that
\eqref{eq:servicetime} measures.

\begin{figure}[t]
  \centering
  \includegraphics[width=\linewidth]{Images/lifecycle.png}
  \caption{Lifecycle of a task. The service time $\zeta_j$ spans both the waiting
  and the active phase, so a controller cannot improve it by delaying assignment.
  The membership sets $\mathcal{Q}_t$, $\mathcal{B}_t$ and $\mathcal{C}_t$ of
  \eqref{eq:lifecycle} are disjoint at every timestep by construction.}
  \label{fig:lifecycle}
\end{figure}

% -----------------------------------------------------------------
\subsection{Who decides}
\label{sec:pf:controllers}

A controller $\pi \in \Pi$ performs two logically distinct jobs. \emph{Task
assignment} maps free agents and waiting tasks to agent--task pairs; \emph{routing}
maps the resulting agent targets and the available system information to actions
$u_i(t) \in \mathcal{U}(\ell_i(t))$. Both jobs are treated as part of $\pi$, so
that a path planner is not described as though it also solves assignment.
Formally, $\Pi$ is the set of pairs % [G3]
$\pi = (\pi^{\mathrm{assign}}, \pi^{\mathrm{route}})$, where
$\pi^{\mathrm{route}}$ is a (possibly stochastic) map from the information
available under the regime in Table~\ref{tab:information} to a joint action in
$\prod_{i=1}^m \mathcal{U}(\ell_i(t))$, and $\pi^{\mathrm{assign}}$ is a map from
free agents and $\mathcal{Q}_t$ to a partial matching between them.
Table~\ref{tab:information} constrains only the argument each map may depend on;
$\Pi$ places no further restriction on how the maps are realised (heuristic,
learned, or otherwise), and the decision problem of
Section~\ref{sec:pf:measures} ranges over all of $\Pi$.

Controllers differ in \emph{what information they may condition on}, and that is
the only distinction the problem formulation needs. Three information regimes are
studied and are summarised in Table~\ref{tab:information}. A centralised
controller may condition on the global state. A section-based controller
partitions the movable vertices into zones,
\begin{equation}
  \mathcal{Y} = \{Y_1,\dots,Y_p\}, \qquad
  \bigcup_{q=1}^{p} Y_q = V_{\mathrm{mov}}, \qquad
  Y_q \cap Y_{q'} = \emptyset \;\; (q \neq q'),
  \label{eq:partition}
\end{equation}
and may condition only on zone-local information, plus an explicit protocol at
zone boundaries. A decentralised controller conditions only on the local
observation defined in Section~\ref{sec:pf:observations} and selects actions from
a shared parameterised policy,
\[
  u_i(t) \sim \pi_\theta\bigl(\cdot \mid o_i(t)\bigr).
\]
Parameter sharing is appropriate for homogeneous agents and does not make
execution centralised; where global information is used during training only, the
method is described as centralised training with decentralised execution.

\begin{table}[t]
  \centering
  \caption{Information regimes. The three architectures are distinguished at the
  problem level only by what each component may condition on. Their algorithmic
  realisations are given in Section~\ref{sec:method:controllers}.}
  \label{tab:information}
  \begin{tabular}{lll}
    \toprule
    Architecture & Assignment conditions on & Routing conditions on \\
    \midrule
    Centralised     & global state & global state \\
    Section-based   & zone-local state, boundary protocol & zone-local state \\
    Decentralised   & global state (held fixed across arms) & $o_i(t)$ only \\
    \bottomrule
  \end{tabular}
\end{table}

Task assignment is held identical across the three architectures. Only routing
differs. This is a design decision of the study rather than a property of the
problem, and it is made so that a difference in results is attributable to the
routing architecture and not confounded with assignment quality.

% -----------------------------------------------------------------
\subsection{What a decentralised agent sees}
\label{sec:pf:observations}

A decentralised agent cannot see the warehouse. It sees a bounded neighbourhood
of its own position, and it can exchange information with nearby agents. Both
bounds are properties of the problem, because they define what a solution is
permitted to depend on.

For observation depth $d$, the visible vertex set and the induced subgraph are
\begin{equation}
  V^{(d)}_i(t) = \{ v \in V_{\mathrm{mov}} : d_G(\ell_i(t), v) \le d \},
  \qquad
  \mathcal{G}^{(d)}_i(t) = G\bigl[V^{(d)}_i(t)\bigr].
  \label{eq:localsubgraph}
\end{equation}
The size of \eqref{eq:localsubgraph} is bounded independently of $|V|$, which is
what allows a single policy to be evaluated on warehouses larger than those it
was trained on. The price is that anything beyond $d$ hops is invisible, so an
agent may commit to a move that becomes a conflict a few steps later.

Direction to a distant target is restored inside the local window by annotating
each visible vertex with its distance to the agent's current target $q_i(t)$,
\begin{equation}
  \eta_i(v,t) = d_G\bigl(v,\, q_i(t)\bigr).
  \label{eq:potential}
\end{equation}
The annotation \eqref{eq:potential} is a potential: by the triangle inequality, % [G-eta]
\begin{equation}
  \eta_i(w,t) - \eta_i(v,t) \;\ge\; -c(v,w) \qquad \text{for every } (v,w) \in E,
  \label{eq:potentialbound}
\end{equation}
with \emph{equality} if and only if $(v,w)$ lies on some shortest path from $v$
to $q_i(t)$. Following an edge that attains equality in
\eqref{eq:potentialbound} realises the shortest path; a merely negative
difference $\eta_i(w,t) - \eta_i(v,t) < 0$ indicates progress toward the target
but does not by itself certify that the move is shortest-path optimal, since the
bound in \eqref{eq:potentialbound} can be strict. Congestion avoidance is
precisely the choice to accept a smaller decrease of $\eta_i$ (or even a
non-negative step) now, in exchange for a lower expected total cost later.
Although $d_G$ may be precomputed, $\eta_i$ itself changes whenever the target
changes, which happens at every reassignment.

A second local quantity, distinct from the potential, summarises how crowded % [G6]
the agent's neighbourhood currently is. For a congestion radius
$r_{\mathrm{cng}}$ --- a parameter separate from the observation depth $d$ and
the communication radius $r_{\mathrm{com}}$ introduced next, and not assumed
equal to either --- and defined so that an agent does not count itself,
\begin{equation}
  \delta_i(t) =
  \frac{\bigl|\{ a_j \in \mathcal{A}\setminus\{a_i\} :
                 \ell_j(t) \in V^{(r_{\mathrm{cng}})}_i(t) \}\bigr|}
       {\max\bigl\{1,\; |V^{(r_{\mathrm{cng}})}_i(t)| - 1\bigr\}} .
  \label{eq:congestion}
\end{equation}

Agents within communication range form a second, time-varying graph whose nodes
are agents rather than locations,
\begin{equation}
  G^{\mathcal{A}}_t = (\mathcal{A}, E^{\mathcal{A}}_t), \qquad
  (a_i, a_j) \in E^{\mathcal{A}}_t \iff
  i \neq j \;\wedge\; d_G(\ell_i(t), \ell_j(t)) \le r_{\mathrm{com}} .
  \label{eq:commgraph}
\end{equation}
The environment graph $G$ and the communication graph $G^{\mathcal{A}}_t$ have
different node sets, different cardinalities and different edge semantics
(``can move to'' versus ``can exchange information with''); they must not be
conflated. The full observation available to agent $a_i$ is
\begin{equation}
  o_i(t) = \bigl( \mathcal{G}^{(d)}_i(t),\; \eta_i(\cdot, t),\; \delta_i(t),\;
                  \text{messages from } \{a_j : (a_i,a_j) \in E^{\mathcal{A}}_t\} \bigr).
  \label{eq:observation}
\end{equation}
How $o_i(t)$ is encoded into a vector, and what a message contains, are modelling
choices specified in Section~\ref{sec:method:model}.
Figure~\ref{fig:observation} illustrates both bounds on one instance.

\begin{figure}[t]
  \centering
  \includegraphics[width=\linewidth]{Images/observation.png}
  \caption{The observation of a single agent. The depth-$d$ induced subgraph is
  drawn at full strength and the rest of the warehouse is dimmed; numbers are the
  target distances $\eta_i(v,t)$. The dashed circle is the communication radius
  $r_{\mathrm{com}}$. One agent lies outside both bounds and is therefore
  invisible, although it may cause a conflict within a few timesteps.}
  \label{fig:observation}
\end{figure}

% -----------------------------------------------------------------
\subsection{The learning problem}
\label{sec:pf:game}

Because each agent acts from $o_i(t)$ while the environment evolves under the
joint action of all agents, the decentralised setting is a partially observable
stochastic game
\begin{equation}
  \mathcal{M} = \bigl( \mathcal{A},\, \mathcal{S},\, \{\mathcal{U}_i\}_{i=1}^{m},\,
                       P,\, \{\mathcal{O}_i\}_{i=1}^{m},\, \{O_i\}_{i=1}^{m},\,
                       \{R_i\}_{i=1}^{m},\, \gamma \bigr),
  \label{eq:posg}
\end{equation}
where $\mathcal{S}$ is the state space, $P$ the transition kernel,
$\mathcal{O}_i$ the observation space, $O_i$ the observation function, and
$R_i : \mathcal{S} \times \mathcal{U}^m \times \mathcal{S} \to \mathbb{R}$ the
reward function. The distinction between spaces and functions on the one hand and
realised states, observations and rewards on the other is maintained throughout;
Appendix~\ref{app:notation} (pointed to from Section~\ref{sec:pf:notation})
disambiguates $\mathcal{O}_i$, $O_i$ and other symbols that share a letter with
an unrelated object elsewhere in this section.
Each agent's action space in $\mathcal{M}$ is % [G-Ui]
$\mathcal{U}_i := \bigcup_{v \in V_{\mathrm{mov}}} \mathcal{U}(v)$, i.e.\ the
union over all vertices of the vertex-local action sets \eqref{eq:actions};
at a state in which $\ell_i(t) = v$, only the subset $\mathcal{U}(v) \subseteq
\mathcal{U}_i$ is enabled; the remaining actions in $\mathcal{U}_i$ have zero
probability under $P$ from that state (equivalently, are masked out), which is
how a graph problem with vertex-dependent action sets is embedded in a POSG with
a fixed per-agent action space.

The congestion feature $\delta_i(t)$ of \eqref{eq:congestion}, included in % [G6]
$o_i(t)$, is available both to the policy and to a reward designer; it is used
when $R_i$ is realised in Section~\ref{sec:method:rl}.
Because the agents are homogeneous and cooperative and share one parameter
vector, the learning objective is stated at team level,
\begin{equation}
  \max_{\theta} J(\theta)
    = \max_{\theta} \mathbb{E}_{\pi_\theta}
      \Bigl[ \sum_{t=0}^{T-1} \gamma^{t}\, \frac{1}{m} \sum_{i=1}^{m} R_i(t) \Bigr].
  \label{eq:objective}
\end{equation}
Equation~\eqref{eq:objective} is written for a parameterised policy $\pi_\theta$; % [G-theta]
it applies unchanged to whichever of the three information regimes of
Table~\ref{tab:information} is realised as a learned controller, with $\theta$
ranging over that regime's own parameters (for instance, the decentralised
regime's shared network weights, or a centralised regime's global policy
parameters). The reward functions $R_i$ are left unspecified here. A reward is an
instrument for obtaining a policy that scores well on the system-level criterion
of Section~\ref{sec:pf:measures}; it can be redesigned without changing the
problem, and it is therefore specified in Section~\ref{sec:method:rl}.

% -----------------------------------------------------------------
\subsection{How the two layers close the loop}
\label{sec:pf:coupling}

Let $s_t \in \mathcal{S}$ contain the agent states, the task lifecycle sets and
the current storage configuration. The system runs on two clocks. On every
timestep, the fast dynamics advance under the controller,
\[
  s_{t+1} \sim P(\cdot \mid s_t, u_t), \qquad u_t \sim \pi(\cdot \mid s_t),
\]
with the corresponding product of local policies under decentralised execution.
At storage epochs only, the slow dynamics additionally apply
\eqref{eq:storageupdate}.

The loop closes in two directions. Forward, a change in $x_t$ changes which
vertices appear as pickups through the coupling condition
\eqref{eq:coupling}, and therefore changes the routing problem. Backward, the
realised routing statistics $\mu_t$ and $w_t$ enter the next storage update
through \eqref{eq:storageupdate}. Neither direction exists in a study that fixes
one layer, and Figure~\ref{fig:loop} shows that the fixed-storage baseline is
exactly this diagram with the backward path removed.

\begin{figure}[t]
  \centering
  \includegraphics[width=\linewidth]{Images/loop.png}
  \caption{The two-timescale loop. The fast path runs every timestep; the slow
  feedback path runs every $\Delta$ steps. Deleting the feedback path yields the
  fixed-storage baseline, so the experimental contrast of this thesis is a
  contrast between this diagram with and without one arrow.}
  \label{fig:loop}
\end{figure}

The separation of timescales has a consequence that constrains the study. Between
epochs, the fast process is approximately stationary and the storage rule sees
only its aggregate statistics. If $\Delta$ is too small relative to the mixing
time of the fast process, $F$ reacts to noise and the task distribution becomes
non-stationary from the learner's point of view; if $\Delta$ is too large, the
estimates $\hat{\rho}, \hat{\mu}, \hat{w}$ are stale. The choice of $\Delta$ is
therefore an experimental variable, not a constant.

% -----------------------------------------------------------------
\subsection{How a run is scored}
\label{sec:pf:measures}

A run is scored on completed work, on the cost of producing it, and on how the
traffic was distributed while producing it. For the completed set $\mathcal{C}_T$
over a horizon $T$, throughput and mean service time are
\begin{equation}
  \Lambda_T = \frac{|\mathcal{C}_T|}{T}, \qquad
  \bar{\zeta}_T = \frac{1}{|\mathcal{C}_T|}
                  \sum_{\tau_j \in \mathcal{C}_T} (d_j - r_j),
  \label{eq:throughput}
\end{equation}
and the movement cost per completed task is
\begin{equation}
  \bar{c}_T = \frac{\sum_{i=1}^{m} \sum_{t=0}^{T-1}
                    \hat{c}\bigl(\ell_i(t), \ell_i(t+1)\bigr)}
                   {|\mathcal{C}_T|} .
  \label{eq:movementcost}
\end{equation}
Runs with $|\mathcal{C}_T| = 0$ are reported separately rather than assigned an
arbitrary ratio. Because a completion-only statistic can improve while the system
falls behind, the backlog $|\mathcal{Q}_T| + |\mathcal{B}_T|$ and the age
distribution of unfinished tasks are always reported alongside
\eqref{eq:throughput}.

Traffic distribution is measured separately from delay. Let $\mu_T(e)$ be the
number of traversals of edge $e$ over the evaluation window,
$E^{+}_T = \{e \in E : \mu_T(e) > 0\}$, and
$p_T(e) = \mu_T(e) / \sum_{e' \in E} \mu_T(e')$. For $|E^{+}_T| > 1$ the
normalised entropy is
\begin{equation}
  H_T = -\frac{1}{\log |E^{+}_T|} \sum_{e \in E^{+}_T} p_T(e) \log p_T(e)
        \;\in [0,1],
  \label{eq:entropy}
\end{equation}
so $H_T = 1$ means perfectly uniform edge usage and $H_T \to 0$ means all traffic
on one edge. By convention, when $|E^{+}_T| = 1$ usage is defined as maximally % [G8]
concentrated, $H_T := 0$ (so $C_T := 1$ below); a run with $|E^{+}_T| = 0$ (no
traversals at all) is reported separately, exactly as for $|\mathcal{C}_T| = 0$
above. Since entropy measures dispersion, the concentration index is its
complement,
\begin{equation}
  C_T = 1 - H_T .
  \label{eq:concentration}
\end{equation}
The direction matters: adding $+H_T$ to a minimised cost would reward funnelling
all agents down one aisle. Two caveats accompany \eqref{eq:entropy}. The
normalisation depends on $|E^{+}_T|$, so runs using very different numbers of
edges are not strictly comparable; and concentration is not congestion, because a
concentrated flow at low load causes no delay. Congestion and concentration are
therefore reported as separate outcomes.
Figure~\ref{fig:concentration} illustrates the two extremes.

\begin{figure}[t]
  \centering
  \includegraphics[width=\linewidth]{Images/concentration.png}
  \caption{Two edge-usage patterns on the same graph, drawn on a shared width
  scale, with the resulting $H_T$ and $C_T$. The panels are schematic
  constructions that define the measure; measured values from experimental runs
  are reported in Section~\ref{sec:experiments}.}
  \label{fig:concentration}
\end{figure}

The system-level criterion aggregates the minimisation-oriented measures,
\begin{equation}
  \mathcal{J}(F,\pi) = q_{\mathrm{svc}}\, \mathbb{E}[\bar{\zeta}_T]
                     + q_{\mathrm{mov}}\, \mathbb{E}[\bar{c}_T]
                     + q_{\mathrm{wait}}\, \mathbb{E}[W_T]
                     + q_{\mathrm{conc}}\, \mathbb{E}[C_T]
                     + q_{\mathrm{back}}\, \mathbb{E}[B_T],
  \label{eq:functional}
\end{equation}
with all $q_\bullet \ge 0$ and $B_T = |\mathcal{Q}_T| + |\mathcal{B}_T|$. $W_T$ is % [G7]
a waiting measure, defined so as not to duplicate $B_T$: rather than queueing
delay before assignment (already captured, in aggregate, by $B_T$'s $|\mathcal{Q}_T|$
term), $W_T$ measures time lost \emph{after} assignment, while a task is
active. For agent $i$ serving task $\tau_j$ over $t \in [y_j, d_j)$, let
$\omega_j(t) \in \{0,1\}$ indicate that $i$ is \emph{blocked} at $t$: it chose
$\mathrm{wait}$, or its chosen move was overridden into a wait by the
conflict-resolution operator of Section~\ref{sec:method:resolution}. Then
\begin{equation}
  W_T = \frac{1}{|\mathcal{C}_T|} \sum_{\tau_j \in \mathcal{C}_T}
        \sum_{t=y_j}^{d_j - 1} \omega_j(t) ,
  \label{eq:waiting}
\end{equation}
using the assignment time $y_j$ of \eqref{eq:lifecycle}; $\omega_j(t)$ is logged
operationally in Section~\ref{sec:impl:logging}. $W_T$ is distinct from the
per-edge estimate $\hat{w}_t(e)$ of \eqref{eq:storageupdate}: the latter is a
forward-looking, causal estimate of downstream contention used only to steer the
storage rule $F$, whereas $W_T$ is a backward-looking, realised summary statistic
used only to score a completed run; the two are never required to agree, and
neither is derived from the other. Throughput is deliberately excluded from
\eqref{eq:functional}: it is maximisation-oriented and, under a fixed arrival
rate, largely redundant with $\bar{\zeta}_T$, so including it with a negative
weight would double-count. It is reported separately.

The formal decision problem is
\begin{equation}
  (F^{\star}, \pi^{\star}) \in
  \arg\min_{F \in \mathcal{F},\, \pi \in \Pi} \mathcal{J}(F,\pi)
  \quad \text{subject to \eqref{eq:feasiblestorage}, \eqref{eq:transition},
  \eqref{eq:vertexconflict}, \eqref{eq:swapconflict} and \eqref{eq:coupling}.}
  \label{eq:decisionproblem}
\end{equation}
Constraints \eqref{eq:vertexconflict}--\eqref{eq:swapconflict} are, by % [G9]
Section~\ref{sec:pf:agents}, enforced by a resolution mechanism applied
identically regardless of which $\pi \in \Pi$ is chosen; under that mechanism
they hold for every $\pi$ and so do not prune $\Pi$ in
\eqref{eq:decisionproblem}. They are listed as constraints to state explicitly
that a solution is a statement about \emph{feasible realised trajectories}, not
because any candidate $(F,\pi)$ could otherwise violate them. Constraints
\eqref{eq:feasiblestorage} and \eqref{eq:coupling}, by contrast, do restrict which
$F$ and which realised task streams are admissible.

Equation~\eqref{eq:decisionproblem} is a framing device. This thesis does not
solve it globally; it estimates the effect of selected $(F,\pi)$ combinations
under a factorial design. Concretely, the empirical question is whether there
exists an adaptive rule $F$ such that
\begin{equation}
  \mathcal{J}(F, \pi) \;<\; \mathcal{J}(F_{\mathrm{fix}}, \pi)
  \quad \text{by a margin exceeding seed variance,}
  \label{eq:rqformal}
\end{equation}
and whether the size of that margin depends on the information regime of $\pi$
given in Table~\ref{tab:information} and on the load parameters $m$ and
$\lambda_{\mathrm{task}}$. Section~\ref{sec:method} realises every symbol
appearing in \eqref{eq:rqformal}.

% -----------------------------------------------------------------
\subsection{Notation used more than once}
\label{sec:pf:notation}

Several symbol families in this section share a base letter across unrelated % [G11]
objects (for instance $c(e)$, $\hat c(v,w)$, $\bar c_T$, $C_T$ and
$\mathcal{C}_t$ are five unrelated quantities). None of these are typos or the
same object under different names. Per the reporting convention of this thesis
--- a full symbol table read on first encounter is a reference manual, not
prose --- the disambiguating table is given once, in
Appendix~\ref{app:notation}, rather than here; consult it if a decorated symbol
in Section~\ref{sec:pf:measures} in particular is unclear.

% -----------------------------------------------------------------
\subsection{Scope and assumptions}
\label{sec:pf:scope}

The model above is deliberately narrower than a physical warehouse, and the
following are not represented. Agents are homogeneous, occupy exactly one vertex,
and carry at most one task; there is no capacity, battery, or kinematic model, so
turning and acceleration costs are absorbed into $c(e)$. Execution is
deterministic: an accepted action always succeeds, so the transition kernel $P$
is stochastic only through order arrivals. The conflict-resolution mechanism of % [G13]
Section~\ref{sec:pf:agents} (which decides who yields) is itself deterministic
given the state, so it introduces no further source of stochasticity. Human
pickers and manual interventions are absent. Inventory relocation is abstract ---
\eqref{eq:storageupdate} changes the placement used to generate future tasks, but
the physical movement of stock is not simulated, and if $\chi = 0$ in
\eqref{eq:storageobjective} that movement is also free. Replenishment and
depletion are modelled only if stated, and otherwise $x_t$ is a placement
template. Tasks carry no deadline: $\tau_j$ fixes only a release time $r_j$, and % [K2]
timeliness is scored only in aggregate, after the fact, through
$\bar{\zeta}_T$ \eqref{eq:throughput}. This is a deliberate narrowing relative to
time-constrained MAPF formulations such as \citet{knippenberg2021}, in which a
per-job deadline is a first-class part of the problem: there, a plan that misses
its deadline is either infeasible (hard constraint) or explicitly penalised in
the objective (soft constraint); here, an arbitrarily late delivery only
degrades $\mathcal{J}(F,\pi)$ through the mean $\bar\zeta_T$, and no individual
task can be said to have failed. Reintroducing deadlines, if the warehouse
motivation calls for service-level guarantees, would mean adding a deadline
$t_j^{\mathrm{d}}$ to \eqref{eq:task} and either a feasibility condition
alongside \eqref{eq:transition}--\eqref{eq:swapconflict} or a term in
\eqref{eq:functional} of the same shape as \citet{knippenberg2021}'s soft
penalty; either is a Problem Formulation-level change, not a Method-level one,
since it would redefine what a valid or well-scored solution is.

Each of these exclusions bounds the claims that the experiments can support. In
particular, results about the value of congestion feedback are conditional on
relocation being cheap, and results about decentralised control are conditional
on deterministic execution.
\end{document}
